% XeLaTeX can use any Mac OS X font. See the setromanfont command below.
% Input to XeLaTeX is full Unicode, so Unicode characters can be typed directly into the source.

% The next lines tell TeXShop to typeset with xelatex, and to open and save the source with Unicode encoding.

%!TEX TS-program = xelatex
%!TEX encoding = UTF-8 Unicode

\documentclass[11pt,twoside,openany]{book}
\usepackage{multirow}
\usepackage{pdfpages}

\usepackage{xspace}

\usepackage{array}
\usepackage{lipsum}

\usepackage{algorithm}
\usepackage{algpseudocode}

\usepackage{geometry}                % See geometry.pdf to learn the layout options. There are lots.
%\usepackage[margin=1cm]{geometry}


% \addtolength{\textwidth}{1.75in}
%
% \addtolength{\topmargin}{-.875in}
% \addtolength{\textheight}{1.75in}

\geometry{a4paper,left=20mm,top=20mm,bottom=50mm,total={162mm,230mm}}                   % ... or a4paper or a5paper or ...
%\geometry{landscape}                % Activate for for rotated page geometry
%\usepackage[parfill]{parskip}    % Activate to begin paragraphs with an empty line rather than an indent
\usepackage{amssymb}
%\usepackage{todonotes}
\setlength{\marginparwidth}{2cm}
\usepackage[backgroundcolor=white,bordercolor=blue,linecolor=blue,textwidth=1cm]{todonotes}
\usepackage{booktabs}
\usepackage{longtable}
% \usepackage{url}
\usepackage{xurl}
\newcommand{\sourcefile}{\url}
\usepackage[depth=4]{bookmark}


\usepackage{algorithm}
\usepackage{algpseudocode}

\usepackage{imakeidx}
\usepackage[utf8]{inputenc}
\usepackage[T1]{fontenc}
\makeindex[intoc]

\usepackage[font=normalsize]{caption}

\renewcommand{\theparagraph}{\S\arabic{paragraph}}
\setcounter{secnumdepth}{4}


\usepackage{xcolor}
\usepackage{listings}

\makeatletter
\global\let\tikz@ensure@dollar@catcode=\relax
\makeatother

\definecolor{mGreen}{rgb}{0,0.6,0}
\definecolor{mGray}{rgb}{0.5,0.5,0.5}
\definecolor{mPurple}{rgb}{0.58,0,0.82}
\definecolor{backgroundColour}{rgb}{0.95,0.95,0.92}
\definecolor{gray}{rgb}{0.4,0.4,0.4}
\definecolor{darkblue}{rgb}{0.0,0.0,0.6}
\definecolor{cyan}{rgb}{0.0,0.6,0.6}

\lstdefinestyle{CStyle}{
    backgroundcolor=\color{backgroundColour},
    commentstyle=\color{mGreen},
    keywordstyle=\color{magenta},
    numberstyle=\tiny\color{mGray},
    stringstyle=\color{mPurple},
    basicstyle=\scriptsize\ttfamily,
    breakatwhitespace=false,
    breaklines=true,
    captionpos=b,
    keepspaces=true,
    numbers=left,
    numbersep=5pt,
    showspaces=false,
    showstringspaces=false,
    showtabs=false,
    tabsize=2,
    language=C
}


\lstset{
  backgroundcolor=\color{white},   % choose the background color; you must add \usepackage{color} or \usepackage{xcolor}; should come as last argument
  basicstyle=\footnotesize\ttfamily,        % the size of the fonts that are used for the code
  breakatwhitespace=false,         % sets if automatic breaks should only happen at whitespace
  breaklines=true,                 % sets automatic line breaking
  captionpos=b,                    % sets the caption-position to bottom
  commentstyle=\color{gray},    % comment style
  deletekeywords={...},            % if you want to delete keywords from the given language
  %escapeinside={\%*}{*)},          % if you want to add LaTeX within your code
  %extendedchars=true,              % lets you use non-ASCII characters; for 8-bits encodings only, does not work with UTF-8
  %firstnumber=1000,                % start line enumeration with line 1000
  frame=single,                    % adds a frame around the code
  keepspaces=true,                 % keeps spaces in text, useful for keeping indentation of code (possibly needs columns=flexible)
  keywordstyle=\color{blue},       % keyword style
  language=C,                 % the language of the code
  %morekeywords={*,...},            % if you want to add more keywords to the set
  numbers=left,                    % where to put the line-numbers; possible values are (none, left, right)
  numbersep=5pt,                   % how far the line-numbers are from the code
  numberstyle=\tiny\color{gray}, % the style that is used for the line-numbers
  rulecolor=\color{black},         % if not set, the frame-color may be changed on line-breaks within not-black text (e.g. comments (green here))
  showspaces=false,                % show spaces everywhere adding particular underscores; it overrides 'showstringspaces'
  showstringspaces=false,          % underline spaces within strings only
  showtabs=false,                  % show tabs within strings adding particular underscores
  stepnumber=1,                    % the step between two line-numbers. If it's 1, each line will be numbered
  stringstyle=\color{black},     % string literal style
  tabsize=2,                     % sets default tabsize to 2 spaces
}

\lstset{
  columns=fullflexible,
  showstringspaces=false,
  commentstyle=\color{gray}\upshape,
  backgroundcolor=\color{backgroundColour},
  commentstyle=\color{mGreen},
  keywordstyle=\color{magenta},
  numberstyle=\tiny\color{mGray},
  stringstyle=\color{mPurple},
  basicstyle=\scriptsize\ttfamily,
  tabsize=2
}

\lstdefinelanguage{XML}
{
  morestring=[b]",
  morestring=[s]{>}{<},
  morecomment=[s]{<?}{?>},
  stringstyle=\color{black},
  identifierstyle=\color{darkblue},
  keywordstyle=\color{cyan},
  morekeywords={xmlns,version,type}% list your attributes here
}


% Will Robertson's fontspec.sty can be used to simplify font choices.
% To experiment, open /Applications/Font Book to examine the fonts provided on Mac OS X,
% and change "Hoefler Text" to any of these choices.

\usepackage{fontspec,xltxtra,xunicode}
\defaultfontfeatures{Mapping=tex-text}
\setromanfont[Mapping=tex-text]{Optima}
\setsansfont[Scale=MatchLowercase,Mapping=tex-text]{Optima}
\setmonofont[Scale=MatchLowercase]{Andale Mono}

%%%% Graphics
\usepackage{graphicx}
\usepackage{tikz}
\definecolor{unilublue}{RGB}{55,149,218}
\definecolor{sntred}{RGB}{219,46,27}
\definecolor{sntpurple}{RGB}{86,30,130}
%\definecolor{sntblue}{RGB}{50,130,207}
\definecolor{sntblue}{RGB}{55,149,218}

\pgfdeclareimage[width=30mm]{logo-snt}{logos/logo-snt}
\pgfdeclareimage[width=30mm]{logo-uni-lu}{logos/logo-uni-lu}

%%%% Fancy
\usepackage{fancyhdr}

%%%% Increase page length
\addtolength{\textheight}{1in}

%%%% Sections
\usepackage{sectsty}
\allsectionsfont{\sffamily}

\setlength{\parskip}{1em}

%\renewcommand{\,}{$^{\cdot}$}

\usepackage{hyperref}
\hypersetup{
    colorlinks,
    citecolor=black,
    filecolor=black,
    linkcolor=black,
    urlcolor=black
}


%%%% Title
\newcommand{\titleone}{\textsf{FAQAS Framework - MASS}}
\newcommand{\titletwo}{\textsf{Verification Report}}
\newcommand{\titlethree}{\textsf{ASN.1 Case Study}}

\newcommand{\todoinline}[1]{\todo[color=orange,inline]{ \textbf{TODO}: #1 }}

\newcommand{\TODO}[1]{\todo[color=orange,inline]{ \textbf{TODO}: #1 }}
\newcommand{\OSCAR}[1]{\todo[color=yellow,inline]{ \textbf{OSCAR}: #1 }}
\newcommand{\DONE}[1]{\todo[color=green,inline]{ \textbf{DONE}: #1 }}


\newcommand{\EMPH}[1]{\textbf{\emph{#1}}}
\newcommand{\INDEX}[1]{\index{\MakeLowercase{#1}}\EMPH{#1}}

\newcommand{\MASS}{\textit{MASS}\xspace}

\newcommand{\SEMUS}{\textit{SEMuS}\xspace}

\newcommand{\SEMU}{\textit{SEMu}\xspace}
\newcommand{\remark}{\textbf{Remark:}\xspace}


\begin{document}

\pagestyle{fancy}
\renewcommand{\sectionmark}[1]{\markright{\textit{#1}}}

\renewcommand{\headrulewidth}{2pt}% 2pt header rule
\renewcommand{\headrule}{\hbox to\headwidth{%
  \color{sntblue}\leaders\hrule height \headrulewidth\hfill}}

\fancyhf{}

%\lhead{\fancyplain{}{\setlength{\unitlength}{1mm}
%\begin{picture}(0,0)
%\put(0,-4){\includegraphics[width=50pt]{logos/logo-snt}}
%\end{picture}}}

\newcommand{\DOCUMENTID}{ITT-1-9873-ESA-FAQAS-SCF}
\lhead{\fancyplain{}{\textit{\DOCUMENTID}}}
\rhead{\fancyplain{}{\rightmark }}

\fancyfoot[C]{%
\begin{tikzpicture}[remember picture,overlay]
\path [fill=sntred]    ([xshift=88pt,yshift=20pt]current page.south west) rectangle
                       ([xshift=229pt,yshift=50pt] current page.south west);
\path [fill=sntpurple] ([xshift=229pt,yshift=20pt] current page.south west) rectangle
                       ([xshift=370pt,yshift=50pt] current page.south west);
\path [fill=sntblue]   ([xshift=370pt,yshift=20pt] current page.south west) rectangle
                       ([xshift=510pt,yshift=50pt] current page.south west);
\end{tikzpicture}
}

\fancyfoot[RO]{
\begin{tikzpicture}[remember picture,overlay]
\node [circle, ultra thick, fill=white, draw=sntblue] at ([xshift=530pt,yshift=35pt] current page.south west) {\thepage};
\end{tikzpicture}}

\fancyfoot[LE]{
\begin{tikzpicture}[remember picture,overlay]
\node [circle, ultra thick, fill=white, draw=sntblue] at ([xshift=65pt,yshift=35pt] current page.south west) {\thepage};
\end{tikzpicture}}

\newcommand{\CHANGED}[1]{{#1}}

\newcommand{\FAQAS}{FAQAS-framework\xspace}


\thispagestyle{empty}

\begin{tikzpicture}[remember picture,overlay]
\path [fill=sntred]    ([xshift=30pt,yshift=20pt]current page.south west) rectangle
                       ([xshift=210pt,yshift=50pt] current page.south west);
\path [fill=sntpurple] ([xshift=210pt,yshift=20pt] current page.south west) rectangle
                       ([xshift=390pt,yshift=50pt] current page.south west);
\path [fill=sntblue]   ([xshift=390pt,yshift=20pt] current page.south west) rectangle
                       ([xshift=570pt,yshift=51pt] current page.south west);
\path [fill=unilublue] ([xshift=30pt,yshift= 50pt] current page.south west) --
                       ([xshift=570pt,yshift= 50pt] current page.south west)
                       [rounded corners=20pt] --
                       ([xshift=570pt,yshift=740pt] current page.south west)
                       [sharp corners] --
                       ([xshift=30pt,yshift=740pt] current page.south west);
\node [fill=white,rounded corners=0pt,inner xsep=6pt,inner ysep=3pt]
      at ([xshift=523pt,yshift=120pt] current page.south west)
      {\pgfuseimage{logo-uni-lu}};

\node [fill=white,rounded corners=2pt,inner xsep=6pt,inner ysep=3pt]
      at ([xshift=520pt,yshift=780pt] current page.south west)
      {\pgfuseimage{logo-snt}};

%\node [circle, fill=white, draw=sntblue] at ([xshift=550pt,yshift=35pt] current page.south west) {a};

\node[draw=none,fill=none,right] at (-1, -7){\color{white}\LARGE\bf\titleone};
\node[draw=none,fill=none,right] at (-1, -8){\color{white}\LARGE\bf\titletwo};
\node[draw=none,fill=none,right] at (-1, -9){\color{white}\LARGE\bf\titlethree};

\node[draw=none,fill=none,right] at (-1, -12){\color{white}\Large\textsf{O. Cornejo, F. Pastore, E. Viganò}};
\node[draw=none,fill=none,right] at (-1, -13){\color{white}\Large\textsf{Interdisciplinary Centre for Security, Reliability and Trust}};
\node[draw=none,fill=none,right] at (-1, -14){\color{white}\Large\textsf{University of Luxembourg}};
\node[draw=none,fill=none,right] at (11, -16){\color{white}\textsf{ITT-1-9873-ESA-FAQAS-SUTP}};
\node[draw=none,fill=none,right] at (11, -17){\color{white}\Large\textsf{Issue 4, Rev. 1}};
\node[draw=none,fill=none,right] at (11, -18){\color{white}\Large\textsf{\today}};
\node[draw=none,fill=none,right] at (-1, -24){\color{white}\tiny\textsf{EUROPEAN SPACE AGENCY. CONTRACT REPORT.}};
\node[draw=none,fill=none,right] at (-1, -24.3){\color{white}\tiny\textsf{The work described in this report was done under ESA contract. Responsibility for the contents resides in the author or organisation that prepared it.}};
\node[draw=none,fill=none,right] at (-1, -24.8){\color{white}\tiny\textsf{The copyright in this document is vested in the University of Luxembourg.}};
\node[draw=none,fill=none,right] at (-1, -25.1){\color{white}\tiny\textsf{This document may only be reproduced in whole or in part, or stored in a retrieval system,or transmitted in any form, or by any means electronic,}};
\node[draw=none,fill=none,right] at (-1, -25.4){\color{white}\tiny\textsf{mechanical, photocopying or otherwise, either with the prior permission of the University of Luxembourg or in accordance with the terms of ESTEC Contract No. 4000128969/19/NL/AS.}};

\node[draw=none,fill=none,right] at (-1, -24){\color{white}\tiny\textsf{EUROPEAN SPACE AGENCY. CONTRACT REPORT.}}; 
\node[draw=none,fill=none,right] at (-1, -24.3){\color{white}\tiny\textsf{The work described in this report was done under ESA contract. Responsibility for the contents resides in the author or organisation that prepared it.}};
\node[draw=none,fill=none,right] at (-1, -24.8){\color{white}\tiny\textsf{The copyright in this document is vested in the University of Luxembourg.}};
\node[draw=none,fill=none,right] at (-1, -25.1){\color{white}\tiny\textsf{This document may only be reproduced in whole or in part, or stored in a retrieval system,or transmitted in any form, or by any means electronic,}}; 
\node[draw=none,fill=none,right] at (-1, -25.4){\color{white}\tiny\textsf{mechanical, photocopying or otherwise, either with the prior permission of the University of Luxembourg or in accordance with the terms of ESTEC Contract No. 4000128969/19/NL/AS.}};


\end{tikzpicture}

\newpage


% !TEX root = MAIN.tex

\section*{Revisions}
\label{sec:revisions}


\setlength\LTleft{0pt}
\setlength\LTright{0pt}
\tiny 
%@{\extracolsep{\fill}}
\begin{longtable}{|p{2cm}|p{1cm}|p{1.5cm}|p{9cm}|@{}}
\label{table:codeoperators} \\
\hline
\textbf{Issue Number}&\textbf{Date}&\textbf{Authors}&\textbf{Description}\\
\hline
ITT-1-9873-ESA-FAQAS-FR
Issue 1 Rev. 1&
October 29th, 2021&
Fabrizio Pastore, Oscar Cornejo, Enrico Viganò&
\begin{minipage}{8cm}
Initial release.
\end{minipage}
\\
\hline
ITT-1-9873-ESA-FAQAS-FR
Issue 1 Rev. 2&
November 11th, 2021&
Fabrizio Pastore, Oscar Cornejo, Enrico Viganò&
\begin{minipage}{8cm}
Added Chapter~\ref{chap:deliverables_summary} (deliverables summary).\\
Added Chapter~\ref{ch:toolset} (description of toolset package).\\
Added information about input and outputs of the tools in Chapter~\ref{chapter:methodology} (see blue text).
\end{minipage}
\\
\hline
                                                    
\end{longtable}
\normalsize

\clearpage


\tableofcontents


% !TEX root = MAIN.tex

\chapter{Scope and content}

This document is the deliverable SSS of the ESA activity ITT-1-9873-ESA. It concerns requirements specification for the \emph{FAQAS framework} to be delivered by ITT-1-9873-ESA. Following the structure described in the SoW \emph{AO9873-ws00pe\_SOW.pdf}, it provides the structured requirements baseline for the FAQAS framework according to ECSS-E-ST-40C Annex B.
 
\section{Applicable and reference documents}

\begin{itemize}
\item{D1 - Mutation testing survey}
\item{D2 - Study of mutation testing applicability to space software}
\end{itemize}

\chapter{Terms, definitions and abbreviated terms}

\begin{itemize}
\item{FAQAS}: activity ITT-1-9873-ESA
\item{FAQAS-framework}: software system to be released at the end of WP4 of FAQAS
\item{D2}: Deliverable D2 of FAQAS, \emph{Study of mutation testing applicability to space software}
\item{KLEE}: Third party test generation tool, details are provided in D2.
\item{SUT}: Software under test, i.e, the software that should be mutated by means of mutation testing.
\item{WP}: Work package
\end{itemize}

\clearpage
 



% !TEX root = MAIN.tex

\chapter{Configuration of the FAQAS Toolset}

% \TODO{that we used}



We configured MASS and SEMuS by following the default configuration provided in the Software User Manual (SUM) document, except for a couple of additional configurations, which are described in this chapter.

\section{MASS Configurations}


Concerning the configuration of the \textbf{MASS tool}, the additional configurations that we considered were the ones regarding the disregarding of equivalent and duplicate mutants with trivial compiler optimizations, and the one concerning the use of sampling.
In the case of trivial compiler optimisations we enabled all levels of optimizations permitted by GCC~\cite{GCCopt}, that is, levels \texttt{O0}, \texttt{O1}, \texttt{O2}, \texttt{O3}, \texttt{Ofast}, and \texttt{Os}. 
In the case of sampling, we decided to analyze all the mutants generated by the tool, and do not employ sampling methods. 

\section{SEMuS Configurations}

Now we report the configurations concerning the \textbf{SEMuS tool}, the additional configurations concerned the compilation command used, the timeouts, and the memory limits imposed to the test generation process.

\begin{lstlisting}[language=bash, label=libutil_com_cmd, caption=Compilation command for a single Libutil source used in SEMuS.]
-g -std=gnu99 -m64 -O0 -Wall -Wextra -Wshadow -Wcast-align -Wwrite-strings -Wno-unused-parameter -I$repo_root_dir/src -I$repo_root_dir/src/gosh -I$repo_root_dir/src/zip -I$repo_root_dir/src/watchdog  -I$repo_root_dir/build/src/watchdog -I$repo_root_dir/build/src -I$repo_root_dir/include -I$repo_root_dir/build/include -I$repo_root_dir/include/gs -I$repo_root_dir/build/include/gs -I$repo_root_dir/include/deprecated -I$repo_root_dir/build/include/deprecated -I$repo_root_dir/include/deprecated/gs/gosh -I$repo_root_dir/build/include/deprecated/gs/gosh -I$repo_root_dir/include -o $out_file $in_file
\end{lstlisting}

Listing~\ref{libutil_com_cmd} presents the compilation command we used to enable test generation on Libutil, in particular, the variable \texttt{\$repo\_root\_dir} represents the installation path of Libutil, \texttt{\$out\_file} represents the file storing the compiled object, and \texttt{\$in\_file} represents the source being analyzed.


We configured SEMuS to terminate the test generation after 300 seconds. The timeout implies that after 300 seconds, SEMuS will try to generate test inputs for the states reached so far, even if the symbolic execution has not reached the end of the function under test. If tests are generated, it means that the processed mutants have been killed. Otherwise, it means that SEMuS could not find test inputs that kill those mutants (they require manual inspection because they might be either equivalent or non-equivalent mutants).

%\DONE{What are the implication of the above? You should write that it still checks if the generated inputs kill the mutants, if not they are discarded.}

We also set SEMuS to use a maximum of 2\,000 MB for the test generation, which means that if the test generation reaches 2\,000 MB of memory consumption, the test generation process will be stopped. The limit of 2\,000 MB was selected by following the guidelines of the underlying test generation tool (i.e., SEMu~\cite{titcheu2021killing}).

%\DONE{What is the motivation for the choice above? Why do we believe is a good compromise between scalability and effectiveness?}



% !TEX root = MAIN.tex

\chapter{Mutation Analysis Results}
\label{mass_output}

% \TODO{Title should be about 'Results' not 'Process'}

%\DONE{Re-read the paper to fix third person singular/plural errors}

In this chapter, we report the outcome of the mutation analysis process through the use of MASS. 
Mutation analysis assess the quality of test suites by injecting artificial faults in the source code (i.e., a mutant), and then measures the number of mutated versions that are detected by the test suite. We report the result in terms of the mutation score, that is, the number of killed mutants (i.e., artificial faults detected by the test suite) divided by the total number of mutants.

% \DONE{This is a verification report, which means that we need to provide action items for the company. What is the suggestions form below, should they reach 100\% code coverage?}



Concerning the components of Libutil to be considered on the mutation analysis process, GomSpace indicated a subset of files from the library, that is not of interest for them, because they are source code belonging to third parties. In consequence, we exclude these sources from the analysis; the sources are listed in Table~\ref{table:not_interesting}. 
Instead, the sources of interest for GomSpace are listed in Table~\ref{table:interesting}.

% !TEX root =  ../MAIN.tex

\begin{table}[h]
\scriptsize
\centering
\caption{Source files not of interest for GomSpace.}
\label{table:not_interesting}
\begin{tabular}{l}
\hline
\textbf{Source file}\\
\hline
src/crc32.c\\
src/linux/drivers/gpio/gpio\_sysfs.c\\
src/linux/drivers/gpio/gpio.c\\
src/linux/drivers/gpio/gpio\_virtual.c\\
src/linux/drivers/i2c/i2c.c\\
src/linux/drivers/i2c/i2c\_native.c\\
src/linux/drivers/sys/memory.c\\
src/linux/drivers/spi/spi.c\\
src/linux/cwd.c\\
src/linux/function.c\\
src/linux/clock.c\\
src/linux/mutex.c\\
src/linux/command\_line.c\\
src/linux/time.c\\
src/linux/thread.c\\
src/linux/signal.c\\
src/linux/queue.c\\
src/linux/sem.c\\
src/linux/delay.c\\
src/linux/stdio.c\\
src/strtoint.c\\
src/base16.c\\
src/log/appender/console.c\\
src/log/appender/simple\_file.c\\
src/log/commands.c\\
src/log/log.c\\
src/crc8.c\\
src/zip/miniz/miniz.c\\
\hline
\end{tabular}
\end{table}

% !TEX root =  ../MAIN.tex

\begin{table}[tb]
\scriptsize
\centering
\caption{Source files of interest for GomSpace.}
\label{table:interesting}
\begin{tabular}{l}
\hline
\textbf{Source file}\\
\hline
src/drivers/sys/memory.c\\
src/byteorder.c\\
src/clock.c\\
src/error.c\\
src/hexdump.c\\
src/vmem/vmem.c\\
src/vmem/commands.c\\
src/zip/zip.c\\
src/timestamp.c\\
src/bytebuffer.c\\
src/function\_scheduler.c\\
src/time.c\\
src/gosh/console.c\\
src/gosh/command.c\\
src/gosh/default\_commands.c\\
src/fletcher.c\\
src/test/cmocka.c\\
src/test/log.c\\
src/test/command.c\\
src/watchdog/monitor\_task.c\\
src/watchdog/watchdog.c\\
src/string.c\\
src/stdio.c\\
\hline
\end{tabular}
\end{table}

%The code-driven mutation analysis process in Libutil targeted all the components covered by the Libutil unit test suite. 
The Libutil unit test suite do not cover hardware-specific functions (e.g., drivers), which are covered at the system level.

\section{Test Suite Code Coverage}

Prior applying mutation analysis to Libutil, we calculated the code coverage of the test suite. Table~\ref{table:gslibutil_coverage} provides the code coverage information of the unit test suite for the Libutil library. Note that the code coverage reported here does not include the source files from Table~\ref{table:not_interesting}. 

\begin{table}[h]
\centering
\scriptsize
\begin{tabular}{|l|l|}
\hline
\textbf{Coverage Type} & \textbf{Coverage Rate} \\
\hline
Statement     & 60.2\% (1\,754 of 2\,915 statements)\\
Functions     & 66.1\% (201 of 304 functions)\\
Branches      & 52.2\% (1\,025 of 1\,965 branches)\\
\hline
\end{tabular}
\caption{Libutil code coverage.}
\label{table:gslibutil_coverage}
\end{table}


\EMPH{Action Item 1}{To improve the quality of the unit test suite, ideally GSL should (1) improve the statement and branch coverage, and make both metrics reach 100\%, or (2) justify the elements not covered by the test suite.}

% \TODO{We know that a subset of the files are third party and not of interest from the test suite; we should exclude them, no?}


\section{MASS Results}

After applying MASS to the software artifact, we report in the following the metrics produced by the tool:

\begin{itemize}
	\item MASS generates a total of 8\,572 mutants, this set of mutants was generated using the operators ABS, AOR, ICR, LCR, ROR, SDL, UOI, AOD, LOD, ROD, BOD, SOD, and LVR.

	\item 2\,575 mutants are filtered by trivial compiler optimisations, that is, mutants that are equivalent with respect to the original program or redundant with respect to other mutants.

	\item In total, MASS analyzed 5\,997 mutants, the process took 962 minutes (16 hours) to complete.

	\item The Libutil test suite identified a total of 4\,440 mutants (i.e., killed mutants).

	\item The Libutil test suite did not identify a total of 1\,557 mutants (i.e., live mutants).

	\item From the set of live mutants, MASS identified additionally 842 likely equivalent mutants based on code coverage.

	\item The MASS mutation score is 86.13\%.

	\item The statement coverage is 60.2\%.
	
	\item The total number of statements covered is 1\,754.
	
	\item The minimum statements covered per source file is 2, and the maximum is 347 statements per source file.
\end{itemize}

\EMPH{Note}{It has been shown that test suites with mutation score greater than 75\% have a high fault detection rate~\cite{chekam2017empirical}. This implies that the Libutil test suite has a high fault detection rate for the statements that are covered. Please note that such observation may not hold for the whole program since code coverage is low (i.e., 60.2\% in this case).}

%\DONE{Recompute the numbers above for the subset of files that are of interest for GSL}

% \TODO{What do we conclude form the above? Are you excluding source files that are not of interest for LXS?}

\subsection{Libutil Live Mutants}

In this section, we provide detailed information about five mutants that present a behavior that largely differs from that of the original software. These mutants shall be discovered by the test suite.

\subsubsection{Live Mutant \#1}

\begin{itemize}
\item \textbf{Source file:} \texttt{src/fletcher.c}
\item \textbf{Mutant ID:} \url{fletcher.mut.65.1_1_9.LOD.gs_fletcher16_update}
\end{itemize}

\begin{lstlisting}[language=C, label=live_1, caption=Original source code.]
void gs_fletcher16_update(gs_fletcher16_t * f16, const void * data_in, size_t size)
{
    if (f16 && data_in && size) {
        const uint8_t * data = data_in;
        for (unsigned int idx = 0; idx < size; ++idx) {
            f16->sum1 = (uint16_t)((f16->sum1 + (*data++)) % 255);
            f16->sum2 = (uint16_t)((f16->sum2 + f16->sum1) % 255);
        }
    }
}
\end{lstlisting}

\begin{lstlisting}[language=C, label=live_1, caption=Mutant source code.]
void gs_fletcher16_update(gs_fletcher16_t * f16, const void * data_in, size_t size)
{
    if (f16) { // first and second operand deleted
        const uint8_t * data = data_in;
        for (unsigned int idx = 0; idx < size; ++idx) {
            f16->sum1 = (uint16_t)((f16->sum1 + (*data++)) % 255);
            f16->sum2 = (uint16_t)((f16->sum2 + f16->sum1) % 255);
        }
    }
}
\end{lstlisting}


\EMPH{Action Item 2}{It seems that the test cases are not verifying a scenario in which both parameters \texttt{f16} and \texttt{data} are null. Maybe GSL has done it through code inspection; however, this may lead to regressions faults not being detected by the test suite in the future (e.g., you modify the code and remove the null check). It is advisable to have such test cases in the unit test suite.}


\subsubsection{Live Mutant \#2}

\begin{itemize}
\item \textbf{Source file:} \texttt{src/watchdog/monitor\_task.c}
\item \textbf{Mutant ID:} \url{monitor_task.mut.57.1_1_11.LOD.gs_swwd_monitor_task_stop}
\end{itemize}

\begin{lstlisting}[language=C, label=live_2, caption=Original source code.]
gs_error_t gs_swwd_monitor_task_stop(uint32_t timeout_s)
{
    gs_swwd_monitor_task_t * monitor = gs_swwd_get_monitor_task();
    GS_CHECK_SUPPORTED(monitor != NULL); /* SWWD must be initialized */
 
    /* Signal the task to stop */
    monitor->do_stop = true;
 
    /* Wait for the task to stop */                                                                                                  
    const uint32_t timeout = GS_SWWD_CHECK_INTERVAL_MS + (timeout_s * 1000);
    uint32_t tm = 0;
    while(monitor->is_running && (tm < timeout)) {
        gs_thread_sleep_ms(100);
        tm += 100;
    }
 
    if (monitor->is_running) {
        return GS_ERROR_BUSY;
    }
 
    return GS_OK;
}
\end{lstlisting}

\begin{lstlisting}[language=C, label=live_2, caption=Mutant source code.]
gs_error_t gs_swwd_monitor_task_stop(uint32_t timeout_s)
{
    gs_swwd_monitor_task_t * monitor = gs_swwd_get_monitor_task();
    GS_CHECK_SUPPORTED(monitor != NULL); /* SWWD must be initialized */
 
    /* Signal the task to stop */
    monitor->do_stop = true;
 
    /* Wait for the task to stop */                                                                                                  
    const uint32_t timeout = GS_SWWD_CHECK_INTERVAL_MS + (timeout_s * 1000);
    uint32_t tm = 0;
    while( (tm < timeout)) { // first while operator deleted
        gs_thread_sleep_ms(100);
        tm += 100;
    }
 
    if (monitor->is_running) {
        return GS_ERROR_BUSY;
    }
 
    return GS_OK;
}
\end{lstlisting}


\EMPH{Action Item 3}{It seems that GSL is not testing the execution of stop when the monitor is not running.}

\subsubsection{Live Mutant \#3}

\begin{itemize}
\item \textbf{Source file:} \texttt{src/clock.c}
\item \textbf{Mutant ID:} \url{clock.mut.88.2_1_31.LCR.gs_clock_from_string}
\end{itemize}

\begin{lstlisting}[language=C, label=live_3, caption=Original source code.]
#if !__AVR__
    // check for ISO8601
    {
        struct tm tm;
        memset(&tm, 0, sizeof(tm)); // no daylight saving
        //int res = sscanf(str, "%" SCNd32 "-%" SCNd32 "-%" SCNd32 "T%" SCNd32 ":%" SCNd32 ":%" SCNd32 "Z",
        int res = sscanf(str, "%d-%d-%dT%d:%d:%dZ",
                         &tm.tm_year, &tm.tm_mon, &tm.tm_mday,
                         &tm.tm_hour, &tm.tm_min, &tm.tm_sec);
        if ((res == 6) &&
            (tm.tm_year >= 1970) && (tm.tm_year <= 2038) &&
            (tm.tm_mon >= 1) && (tm.tm_mon <= 12) &&
            (tm.tm_mday >= 1) && (tm.tm_mday <= 31) &&
            (tm.tm_hour >= 0) && (tm.tm_hour <= 23) &&
            (tm.tm_min >= 0) && (tm.tm_min <= 59) &&
            (tm.tm_sec >= 0) && (tm.tm_sec <= 60))
        {
            tm.tm_year -= 1900;
            tm.tm_mon -= 1;
 
#if __linux__
            // not posix compliant
            time_t sec = timegm(&tm);
#else
            // embedded platforms do not have timezones/daylight-saving - so standard mktime works
            time_t sec = mktime(&tm);
#endif
            if (sec >= 0) {
                ts->tv_sec = (uint32_t) sec;
                ts->tv_nsec = 0;
                return GS_OK;
            }
        }
    }
#endif

\end{lstlisting}

\begin{lstlisting}[language=C, label=live_3, caption=Mutant source code.]
#if !__AVR__
    // check for ISO8601
    {
        struct tm tm;
        memset(&tm, 0, sizeof(tm)); // no daylight saving
        //int res = sscanf(str, "%" SCNd32 "-%" SCNd32 "-%" SCNd32 "T%" SCNd32 ":%" SCNd32 ":%" SCNd32 "Z",
        int res = sscanf(str, "%d-%d-%dT%d:%d:%dZ",
                         &tm.tm_year, &tm.tm_mon, &tm.tm_mday,
                         &tm.tm_hour, &tm.tm_min, &tm.tm_sec);
        if ((res == 6) &&
            (tm.tm_year >= 1970) && (tm.tm_year <= 2038) &&
            (tm.tm_mon >= 1) && (tm.tm_mon <= 12) &&
            (tm.tm_mday >= 1) || (tm.tm_mday <= 31) && // '&&'' replaced by a '||'
            (tm.tm_hour >= 0) && (tm.tm_hour <= 23) &&
            (tm.tm_min >= 0) && (tm.tm_min <= 59) &&
            (tm.tm_sec >= 0) && (tm.tm_sec <= 60))
        {
            tm.tm_year -= 1900;
            tm.tm_mon -= 1;
 
#if __linux__
            // not posix compliant
            time_t sec = timegm(&tm);
#else
            // embedded platforms do not have timezones/daylight-saving - so standard mktime works
            time_t sec = mktime(&tm);
#endif
            if (sec >= 0) {
                ts->tv_sec = (uint32_t) sec;
                ts->tv_nsec = 0;
                return GS_OK;
            }
        }
    }
#endif
\end{lstlisting}


\EMPH{Action Item 4}{This mutant shows that GSL does not verify the result returned by the function when the \texttt{day} field of the date contains a value higher than 31.}


\subsubsection{Live Mutant \#4}

\begin{itemize}
\item \textbf{Source file:} \texttt{src/error.c}
\item \textbf{Mutant ID:} \url{error.mut.61.1_1_5.SDL.gs_error_string}
\end{itemize}

\begin{lstlisting}[language=C, label=live_4, caption=Original source code.]
const char * gs_error_string(int error)
{
    switch (error) {
        case GS_OK: return "GS_OK(0)";
        case GS_ERROR_PERM: return GS_DEF2STRING(GS_ERROR_PERM) "(-1)";
        case GS_ERROR_INTR: return GS_DEF2STRING(GS_ERROR_INTR) "(-4)";
            ...
        case GS_ERROR_AMBIGUOUS: return GS_DEF2STRING(GS_ERROR_AMBIGUOUS) "(-2009)";
        case GS_ERROR_STATE: return GS_DEF2STRING(GS_ERROR_STATE) "(-2010)";
    }
    
    // as fallback we use standard POSIX error string
    const int posix_error = abs(error);
    return strerror(posix_error);
}
\end{lstlisting}

\begin{lstlisting}[language=C, label=live_4, caption=Mutant source code.]
const char * gs_error_string(int error)
{
    switch (error) {
        case GS_OK: return "GS_OK(0)";
        case GS_ERROR_PERM: return GS_DEF2STRING(GS_ERROR_PERM) "(-1)";
        case GS_ERROR_INTR: return GS_DEF2STRING(GS_ERROR_INTR) "(-4)";
            ...
        case GS_ERROR_AMBIGUOUS: return GS_DEF2STRING(GS_ERROR_AMBIGUOUS) "(-2009)";
        case GS_ERROR_STATE: return GS_DEF2STRING(GS_ERROR_STATE) "(-2010)";
    }
   
    // as fallback we use standard POSIX error string
    const int posix_error = abs(error);
    ; // return statement deleted
}
\end{lstlisting}


\EMPH{Action Item 5}{This mutant shows that the return string of the function \texttt{gs\_error\_string} is never checked, maybe because it is used only in logging functions.}

\subsubsection{Live Mutant \#5}

\begin{itemize}
\item \textbf{Source file:} \texttt{src/string.c}
\item \textbf{Mutant ID:} \url{string.mut.151.2_1_28.LCR.gs_string_to_uint32}
\end{itemize}

\begin{lstlisting}[language=C, label=live_5, caption=Original source code.]
 // check for hexadecimal notation
if ((string[0] == '0') && ((string[1] == 'x') || (string[1] == 'X')))                                                             
{
    base = 16;
}
\end{lstlisting}

\begin{lstlisting}[language=C, label=live_5, caption=Mutant source code.]
// check for hexadecimal notation
if ((string[0] == '0') || ((string[1] == 'x') || (string[1] == 'X')))                                                              
{
    base = 16;
}
\end{lstlisting}


\EMPH{Action Item 6}{To cover the mutant, the solution would be to add the following assertion in the test case \texttt{TEST\_gs\_log\_string\_to\_mask}: 
\begin{lstlisting}[language=C, label=assertion]
GS_ASSERT_ERROR_EQUAL(gs_log_string_to_mask("012", 0xff, &new_mask), GS_ERROR_DATA);
\end{lstlisting}

However, such assertion would fail, because the original GSL code returns a \texttt{GS\_OK} code. This might be a fault in the implementation.
}




% \DONE{Report the mutants that we have already shared with them. Avoid duplicate ones.}


% !TEX root = MAIN.tex

\chapter{Test Suite Augmentation Process}

In this chapter, we report the outcome of the test suite augmentation process based on mutation testing via SEMuS. The purpose of this process is to increment the quality of the existing test suite, by adding test cases that kill the mutants not identified in the mutation analysis process (see Chapter~\ref{mass_output}). 
In a nutshell, SEMuS works by processing the SUT source code, together with the list of live mutants, and generates test inputs with the support of a symbolic execution search engine (i.e., KLEE~\cite{cadar2008klee}).


Concerning Libutil, we first identify the source code to be analyzed. We perform this operation for two reasons, (1) because GomSpace defined a subset of source files that were not interesting for them, and (2) because there is a subset of source files that is not compatible with SEMuS (e.g., source files processing I/O operations, source files with external dependencies).  

\TODO{The list not of interest should go to the previous chapter}

In the following, we list the files that were not of interest for GomSpace:

\begin{itemize}
\item src/crc32.c
\item src/linux/drivers/gpio/gpio\_sysfs.c
\item src/linux/drivers/gpio/gpio.c
\item src/linux/drivers/gpio/gpio\_virtual.c
\item src/linux/drivers/i2c/i2c.c
\item src/linux/drivers/i2c/i2c\_native.c
\item src/linux/drivers/sys/memory.c
\item src/linux/drivers/spi/spi.c
\item src/linux/cwd.c
\item src/linux/function.c
\item src/linux/clock.c
\item src/linux/mutex.c
\item src/linux/command\_line.c
\item src/linux/time.c
\item src/linux/thread.c
\item src/linux/signal.c
\item src/linux/queue.c
\item src/linux/sem.c
\item src/linux/delay.c
\item src/linux/stdio.c
\item src/strtoint.c
\item src/base16.c
\item src/log/appender/console.c
\item src/log/appender/simple\_file.c
\item src/log/commands.c
\item src/log/log.c
\item src/crc8.c
\item src/zip/miniz/miniz.c
\end{itemize}

Now, we list the files that could not be processed because of incompatibilities with the tool:

\begin{itemize}
\item src/hexdump.c
\item src/vmem/vmem.c
\item src/vmem/commands.c
\item src/zip/zip.c
\item src/bytebuffer.c
\item src/function\_scheduler.c
\item src/time.c
\item src/gosh/console.c
\item src/gosh/command.c
\item src/gosh/default\_commands.c
\item src/fletcher.c
\item src/test/cmocka.c
\item src/test/log.c
\item src/test/command.c
\item src/watchdog/monitor\_task.c
\item src/watchdog/watchdog.c
\item src/string.c
\item src/stdio.c
\item src/timestamp.c
\end{itemize}

Finally, we list the source files processed by SEMuS:

\begin{itemize}
\item src/drivers/sys/memory.c
\item src/byteorder.c
\item src/clock.c
\item src/error.c
\end{itemize}

\TODO{Rephase the sentence below}
Within the list of live mutants produced by MASS, this four sources are equivalent to 109 mutants, which were correctly processed by SEMuS.

\section{SEMuS Output}

After applying SEMuS to the artifact, we report in the following the metrics produced by the tool:

\begin{itemize}
	\item The total number of analyzed mutants is 109.
	\item The total number of killed mutants is 30.
	\item The total number of live mutants is 79.
	\item There were not timeouts during the test generation process.
	\item \textbf{The final mutation score of Libutil is 71.44\%.}
\end{itemize}

\TODO{What is the 'final mutation score'? We should define how it is computed.}

Furthermore, we report the status of each processed mutant:

\begin{lstlisting}[language=bash, label=statuses, caption=Mutants statuses.]
clock.mut.37.4_3_29.UOI.gs_clock_to_iso8601_string2.c;LIVE
clock.mut.40.1_1_9.LOD.gs_clock_to_iso8601_string2.c;LIVE
clock.mut.40.2_4_10.UOI.gs_clock_to_iso8601_string2.c;LIVE
clock.mut.40.1_2_10.UOI.gs_clock_to_iso8601_string2.c;LIVE
clock.mut.40.4_2_14.ROR.gs_clock_to_iso8601_string2.c;LIVE
clock.mut.40.5_3_14.ROR.gs_clock_to_iso8601_string2.c;LIVE
clock.mut.40.2_1_16.ICR.gs_clock_to_iso8601_string2.c;LIVE
clock.mut.40.1_10_35.ROR.gs_clock_to_iso8601_string2.c;LIVE
clock.mut.40.1_6_31.UOI.gs_clock_to_iso8601_string2.c;LIVE
clock.mut.40.3_9_38.UOI.gs_clock_to_iso8601_string2.c;LIVE
clock.mut.40.4_11_38.UOI.gs_clock_to_iso8601_string2.c;LIVE
clock.mut.41.1_1_9.SDL.gs_clock_to_iso8601_string2.c;LIVE
clock.mut.41.4_4_22.AOR.gs_clock_to_iso8601_string2.c;LIVE
clock.mut.41.2_1_24.ICR.gs_clock_to_iso8601_string2.c;LIVE
clock.mut.60.2_3_14.LOD.gs_clock_from_string.c;LIVE
clock.mut.60.1_1_9.LOD.gs_clock_from_string.c;KILLED
clock.mut.69.2_5_17.ROR.gs_clock_from_string.c;LIVE
clock.mut.80.1_1_9.SDL.gs_clock_from_string.c;LIVE
clock.mut.80.1_2_21.ICR.gs_clock_from_string.c;LIVE
clock.mut.80.2_1_21.ICR.gs_clock_from_string.c;LIVE
clock.mut.85.2_1_24.LCR.gs_clock_from_string.c;LIVE
clock.mut.85.1_11_13.LOD.gs_clock_from_string.c;LIVE
clock.mut.86.2_1_34.LOD.gs_clock_from_string.c;LIVE
clock.mut.86.2_1_34.LCR.gs_clock_from_string.c;LIVE
clock.mut.85.1_10_13.LOD.gs_clock_from_string.c;LIVE
clock.mut.86.2_2_58.LCR.gs_clock_from_string.c;LIVE
clock.mut.85.1_9_13.LOD.gs_clock_from_string.c;LIVE
clock.mut.87.2_1_30.LCR.gs_clock_from_string.c;LIVE
clock.mut.87.2_2_51.LCR.gs_clock_from_string.c;LIVE
clock.mut.85.1_8_13.LOD.gs_clock_from_string.c;LIVE
clock.mut.88.2_1_31.LCR.gs_clock_from_string.c;LIVE
clock.mut.85.1_7_13.LOD.gs_clock_from_string.c;LIVE
clock.mut.88.2_2_53.LCR.gs_clock_from_string.c;LIVE
clock.mut.88.2_2_53.LOD.gs_clock_from_string.c;LIVE
clock.mut.85.1_6_13.LOD.gs_clock_from_string.c;LIVE
clock.mut.89.2_1_31.LCR.gs_clock_from_string.c;KILLED
clock.mut.85.1_5_13.LOD.gs_clock_from_string.c;LIVE
clock.mut.89.2_2_53.LCR.gs_clock_from_string.c;LIVE
clock.mut.85.1_4_13.LOD.gs_clock_from_string.c;LIVE
clock.mut.89.2_2_53.LOD.gs_clock_from_string.c;LIVE
clock.mut.85.1_3_13.LOD.gs_clock_from_string.c;LIVE
clock.mut.90.2_1_30.LCR.gs_clock_from_string.c;LIVE
clock.mut.90.2_1_30.LOD.gs_clock_from_string.c;LIVE
clock.mut.90.2_2_51.LCR.gs_clock_from_string.c;LIVE
clock.mut.90.2_2_51.LOD.gs_clock_from_string.c;LIVE
clock.mut.85.1_2_13.LOD.gs_clock_from_string.c;LIVE
clock.mut.85.1_1_13.LOD.gs_clock_from_string.c;LIVE
clock.mut.91.2_1_30.LCR.gs_clock_from_string.c;LIVE
clock.mut.85.2_2_18.ROD.gs_clock_from_string.c;LIVE
clock.mut.85.4_3_18.ROR.gs_clock_from_string.c;LIVE
clock.mut.85.2_5_18.ROR.gs_clock_from_string.c;LIVE
clock.mut.86.1_1_14.ROD.gs_clock_from_string.c;LIVE
clock.mut.86.2_2_25.ROD.gs_clock_from_string.c;LIVE
clock.mut.86.3_5_28.ICR.gs_clock_from_string.c;LIVE
clock.mut.86.6_4_28.ICR.gs_clock_from_string.c;LIVE
clock.mut.86.5_3_28.ICR.gs_clock_from_string.c;LIVE
clock.mut.86.2_1_28.ICR.gs_clock_from_string.c;LIVE
clock.mut.86.1_6_28.ICR.gs_clock_from_string.c;LIVE
clock.mut.86.4_8_52.ICR.gs_clock_from_string.c;LIVE
clock.mut.87.2_2_24.ROD.gs_clock_from_string.c;LIVE
clock.mut.87.2_1_27.ICR.gs_clock_from_string.c;LIVE
clock.mut.87.3_2_27.ICR.gs_clock_from_string.c;LIVE
clock.mut.87.2_4_44.ROD.gs_clock_from_string.c;LIVE
clock.mut.87.4_5_47.ICR.gs_clock_from_string.c;LIVE
clock.mut.88.1_1_14.ROD.gs_clock_from_string.c;LIVE
clock.mut.88.2_2_25.ROD.gs_clock_from_string.c;LIVE
clock.mut.88.3_2_28.ICR.gs_clock_from_string.c;LIVE
clock.mut.88.2_1_28.ICR.gs_clock_from_string.c;LIVE
clock.mut.88.2_4_46.ROD.gs_clock_from_string.c;LIVE
clock.mut.88.4_5_49.ICR.gs_clock_from_string.c;LIVE
clock.mut.89.2_1_28.ICR.gs_clock_from_string.c;LIVE
clock.mut.89.4_4_49.ICR.gs_clock_from_string.c;LIVE
clock.mut.90.2_1_27.ICR.gs_clock_from_string.c;LIVE
clock.mut.90.4_4_47.ICR.gs_clock_from_string.c;LIVE
clock.mut.91.2_1_27.ICR.gs_clock_from_string.c;LIVE
clock.mut.91.6_6_44.ROR.gs_clock_from_string.c;LIVE
clock.mut.91.4_4_47.ICR.gs_clock_from_string.c;LIVE
clock.mut.93.10_1_24.AOR.gs_clock_from_string.c;LIVE
clock.mut.103.2_1_24.ICR.gs_clock_from_string.c;LIVE
memory.mut.8.3_1_9.UOI.long_to_string.c;KILLED
memory.mut.8.1_1_9.ABS.long_to_string.c;KILLED
memory.mut.8.1_2_9.UOI.long_to_string.c;KILLED
memory.mut.8.4_3_9.UOI.long_to_string.c;KILLED
memory.mut.8.2_4_9.UOI.long_to_string.c;KILLED
memory.mut.8.6_1_16.ROR.long_to_string.c;KILLED
memory.mut.8.4_3_16.ROR.long_to_string.c;KILLED
memory.mut.8.1_1_9.ROD.long_to_string.c;KILLED
memory.mut.8.5_4_16.ROR.long_to_string.c;KILLED
memory.mut.8.1_2_19.ICR.long_to_string.c;KILLED
memory.mut.8.2_1_19.ICR.long_to_string.c;KILLED
memory.mut.9.1_2_23.UOI.long_to_string.c;KILLED
memory.mut.9.3_1_23.UOI.long_to_string.c;KILLED
memory.mut.9.1_1_40.ABS.long_to_string.c;KILLED
memory.mut.9.1_6_40.UOI.long_to_string.c;KILLED
memory.mut.9.3_5_40.UOI.long_to_string.c;KILLED
memory.mut.11.1_1_9.SDL.long_to_string.c;KILLED
memory.mut.11.1_2_23.UOI.long_to_string.c;KILLED
memory.mut.11.3_1_23.UOI.long_to_string.c;KILLED
memory.mut.11.1_1_49.ABS.long_to_string.c;KILLED
memory.mut.11.1_6_49.UOI.long_to_string.c;KILLED
memory.mut.11.3_5_49.UOI.long_to_string.c;KILLED
timestamp.mut.11.1_5_38.ROR.timestamp_add.c;KILLED
timestamp.mut.13.2_7_56.AOR.timestamp_add.c;KILLED
timestamp.mut.26.10_1_18.AOR.timestamp_diff.c;KILLED
timestamp.mut.27.1_5_23.ROR.timestamp_diff.c;KILLED
timestamp.mut.44.2_5_23.ROR.timestamp_ge.c;LIVE
timestamp.mut.44.1_1_10.ROD.timestamp_ge.c;LIVE
timestamp.mut.44.2_2_23.ROD.timestamp_ge.c;KILLED
timestamp.mut.44.4_3_23.ROR.timestamp_ge.c;KILLED
\end{lstlisting}

In the following sections, we provide the list of test cases that covers the 30 mutants killed for the sources \texttt{clock.c}, \texttt{timestamp.c}, \texttt{memory.c}, and \texttt{error.c}.

\TODO{I would report the in an appendix. ALso, we should report both the test case and the oracles.}
\TODO{You should explicitely indicate that engineers are supposed to verify if teh expected output is correct, if not, it means that the software is faulty}.
\TODO{Did you check if the same input kills other mutants? We shall do it, so that we minimize the number of test cases reported}
\TODO{What about the mutants for which SEMuS do not generate inputs but terminate? I would suggest to consider those mutants as likely equivalent and list them as such}
\TODO{The mutants are not many, can we try to analyze them to draw conclusions about the quality of the test suite? For example, can you printout/grep the mutants for which we generated inputs so that we can understand if they are about boundary conditions? Does the nature of the mutation operator help us in this regard? We need to provide an evaluation of the test suite.}

\newpage
\subsection{Test Cases Generated for clock.c}

In the following, we present the test cases that have been generated for the function \texttt{gs\_clock\_from\_string}.

\lstinputlisting{generated_test_cases/direct-src.clock/gs_clock_from_string/FAQAS_SEMu-out/produced-unittests/test000001.ktest.c}

\lstinputlisting{generated_test_cases/direct-src.clock/gs_clock_from_string/FAQAS_SEMu-out/produced-unittests/test000008.ktest.c}

\newpage
\subsection{Test Cases Generated for memory.c}

In the following, we present the test cases that has been generated for the function \texttt{long\_to\_string}.

\lstinputlisting{generated_test_cases/direct-src.drivers.sys.memory/long_to_string/FAQAS_SEMu-out/produced-unittests/test000001.ktest.c}
\lstinputlisting{generated_test_cases/direct-src.drivers.sys.memory/long_to_string/FAQAS_SEMu-out/produced-unittests/test000003.ktest.c}
\lstinputlisting{generated_test_cases/direct-src.drivers.sys.memory/long_to_string/FAQAS_SEMu-out/produced-unittests/test000005.ktest.c}
\lstinputlisting{generated_test_cases/direct-src.drivers.sys.memory/long_to_string/FAQAS_SEMu-out/produced-unittests/test000006.ktest.c}
\lstinputlisting{generated_test_cases/direct-src.drivers.sys.memory/long_to_string/FAQAS_SEMu-out/produced-unittests/test000010.ktest.c}
\lstinputlisting{generated_test_cases/direct-src.drivers.sys.memory/long_to_string/FAQAS_SEMu-out/produced-unittests/test000016.ktest.c}
\lstinputlisting{generated_test_cases/direct-src.drivers.sys.memory/long_to_string/FAQAS_SEMu-out/produced-unittests/test000020.ktest.c}
\lstinputlisting{generated_test_cases/direct-src.drivers.sys.memory/long_to_string/FAQAS_SEMu-out/produced-unittests/test000039.ktest.c}
\lstinputlisting{generated_test_cases/direct-src.drivers.sys.memory/long_to_string/FAQAS_SEMu-out/produced-unittests/test000040.ktest.c}
\lstinputlisting{generated_test_cases/direct-src.drivers.sys.memory/long_to_string/FAQAS_SEMu-out/produced-unittests/test000045.ktest.c}
\lstinputlisting{generated_test_cases/direct-src.drivers.sys.memory/long_to_string/FAQAS_SEMu-out/produced-unittests/test000048.ktest.c}
\lstinputlisting{generated_test_cases/direct-src.drivers.sys.memory/long_to_string/FAQAS_SEMu-out/produced-unittests/test000054.ktest.c}
\lstinputlisting{generated_test_cases/direct-src.drivers.sys.memory/long_to_string/FAQAS_SEMu-out/produced-unittests/test000064.ktest.c}
\lstinputlisting{generated_test_cases/direct-src.drivers.sys.memory/long_to_string/FAQAS_SEMu-out/produced-unittests/test000065.ktest.c}
\lstinputlisting{generated_test_cases/direct-src.drivers.sys.memory/long_to_string/FAQAS_SEMu-out/produced-unittests/test000070.ktest.c}
\lstinputlisting{generated_test_cases/direct-src.drivers.sys.memory/long_to_string/FAQAS_SEMu-out/produced-unittests/test000074.ktest.c}
\lstinputlisting{generated_test_cases/direct-src.drivers.sys.memory/long_to_string/FAQAS_SEMu-out/produced-unittests/test000078.ktest.c}
\lstinputlisting{generated_test_cases/direct-src.drivers.sys.memory/long_to_string/FAQAS_SEMu-out/produced-unittests/test000080.ktest.c}
\lstinputlisting{generated_test_cases/direct-src.drivers.sys.memory/long_to_string/FAQAS_SEMu-out/produced-unittests/test000084.ktest.c}
\lstinputlisting{generated_test_cases/direct-src.drivers.sys.memory/long_to_string/FAQAS_SEMu-out/produced-unittests/test000087.ktest.c}
\lstinputlisting{generated_test_cases/direct-src.drivers.sys.memory/long_to_string/FAQAS_SEMu-out/produced-unittests/test000091.ktest.c}
\lstinputlisting{generated_test_cases/direct-src.drivers.sys.memory/long_to_string/FAQAS_SEMu-out/produced-unittests/test000093.ktest.c}


\newpage
\subsection{Test Cases Generated for error.c}

In the following, we present the test cases that has been generated for the function \texttt{gs\_error\_string}.

\lstinputlisting{generated_test_cases/direct-src.error/gs_error_string/FAQAS_SEMu-out/produced-unittests/test000001.ktest.c}
\lstinputlisting{generated_test_cases/direct-src.error/gs_error_string/FAQAS_SEMu-out/produced-unittests/test000006.ktest.c}
\lstinputlisting{generated_test_cases/direct-src.error/gs_error_string/FAQAS_SEMu-out/produced-unittests/test000009.ktest.c}
\lstinputlisting{generated_test_cases/direct-src.error/gs_error_string/FAQAS_SEMu-out/produced-unittests/test000010.ktest.c}


\newpage
\subsection{Test Cases Generated for timestamp.c}

In the following, we present the test cases that has been generated for the function \texttt{timestamp\_add}.

\lstinputlisting{generated_test_cases/direct-src.timestamp/timestamp_add/FAQAS_SEMu-out/produced-unittests/test000001.ktest.c}
\lstinputlisting{generated_test_cases/direct-src.timestamp/timestamp_add/FAQAS_SEMu-out/produced-unittests/test000002.ktest.c}

In the following, we present the test cases that has been generated for the function \texttt{timestamp\_diff}.

\lstinputlisting{generated_test_cases/direct-src.timestamp/timestamp_diff/FAQAS_SEMu-out/produced-unittests/test000001.ktest.c}
\lstinputlisting{generated_test_cases/direct-src.timestamp/timestamp_diff/FAQAS_SEMu-out/produced-unittests/test000003.ktest.c}


In the following, we present the test cases that has been generated for the function \texttt{timestamp\_ge}.

\lstinputlisting{generated_test_cases/direct-src.timestamp/timestamp_ge/FAQAS_SEMu-out/produced-unittests/test000001.ktest.c}
\lstinputlisting{generated_test_cases/direct-src.timestamp/timestamp_ge/FAQAS_SEMu-out/produced-unittests/test000003.ktest.c}

% !TEX root = MAIN.tex

\chapter{Software Defects}

Listing~\ref{semus:testUtil}, shows the test case generated for a mutant of function \emph{timestamp\_diff} in Libutil.
This function receives as input two parameters, \emph{base} and \emph{diff}; it subtracts \emph{diff} from \emph{base} and updates the value of \emph{base}. Both \emph{base} and \emph{diff} are of type \emph{struct gs\_timestamp\_t}, which is shown in Figure~\ref{semus:testUtil:struct}.


The test case initializes \emph{base} and \emph{diff} as follows:
\begin{verbatim}
base = {tv_sec= 0,tv_nsec= 0};
diff = {tv_sec= 0,tv_nsec=1}
\end{verbatim}

The output generated by the SUT (i.e., the output stored by SEMuS in the file \emph{.expected}) is the following:
\begin{verbatim}
result_faqas_semu = 0
base->tv_sec = 4294967295
base->tv_nsec = 999999999
\end{verbatim}

The output above, which simply reflects the actual behavior of the implementation, show that the difference between zero (i.e, \emph{base = {tv\_sec= 0,tv\_nsec= 0};}) and one nanosecond  (i.e., \emph{diff = {tv\_sec= 0,tv\_nsec=1}})leads to a big integer (i.e., we observe an \EMPH{integer underflow}); instead, according to specifications, it should make the function return an error code (i.e., a value different than zero shall be assigned to \emph{result\_faqas\_semu} in the output above).


\begin{lstlisting}[language=C, caption=Test case generated for Libutil., label=semus:testUtil]
int main(int argc, char** argv)
{
    (void)argc;
    (void)argv;

    // Declare variable to hold function returned value
    int result_faqas_semu;

    // Declare arguments and make input ones symbolic
    gs_timestamp_t base;
    gs_timestamp_t diff;
    memset(&base, 0, sizeof(base));
    memset(&diff, 0, sizeof(diff));
    const unsigned char base_faqas_semu_test_data[] = {0x00, 0x00, 0x00, 0x00, 0x00, 0x00, 0x00, 0x00};
    const unsigned char diff_faqas_semu_test_data[] = {0x00, 0x00, 0x00, 0x00, 0x01, 0x00, 0x00, 0x00};
    memcpy(&base, base_faqas_semu_test_data, sizeof(base));
    memcpy(&diff, diff_faqas_semu_test_data, sizeof(diff));

    // Call function under test
    result_faqas_semu = timestamp_diff(&base, &diff);

    // Make some output
    printf("FAQAS-SEMU-TEST_OUTPUT: result_faqas_semu = %d\n", result_faqas_semu);
    printf("FAQAS-SEMU-TEST_OUTPUT: base->tv_sec = %d\n", base->tv_sec);
    printf("FAQAS-SEMU-TEST_OUTPUT: base->tv_nsec = %d\n", base->tv_nsec);
    return (int)result_faqas_semu;
}
\end{lstlisting}

\begin{lstlisting}[language=C, caption=Definition of \emph{gs\_timestamp\_t} in Libutil., label=semus:testUtil:struct]
struct gs_timestamp_t {
    uint_32t  tv_sec;
    uint_32t  tv_nsec;
}
\end{lstlisting}



\clearpage

\printindex

\end{document}
